Eine Bewegung ist gleichförmig, wenn die Geschwindigkeit konstant ist. Im $s(t)$-Diagramm ist ihr Graph eine Gerade (die Steigung ist die Geschwindigkeit), im $v(t)$-Diagramm eine waagrechte Linie.
\begin{abcliste}
\abc $\result{\text{gleichförmig}}$. Der Ort nimmt linear ab: Der Körper bewegt sich mit konstanter Geschwindigkeit in die negative Richtung, bis er beim Nullpunkt ankommt.
\abc $\result{\text{nicht gleichförmig}}$. Die Geschwindigkeit nimmt gleichmässig zu: Der Körper startet aus dem Stillstand und wird gleichmässig schneller (gleichmässig beschleunigte Bewegung).
\abc $\result{\text{nicht gleichförmig}}$. Die Steigung des Graphen nimmt zu: Der Körper startet ein Stück vom Nullpunkt entfernt mit kleiner Geschwindigkeit und wird immer schneller.
\abc Der Ort ändert sich nicht, der Körper ist in Ruhe. Das ist der Grenzfall einer $\result{\text{gleichförmigen}}$ Bewegung mit der Geschwindigkeit null.
\abc $\result{\text{gleichförmig}}$. Die Geschwindigkeit ist konstant.
\abc $\result{\text{nicht gleichförmig}}$. Der Graph besteht aus zwei Geradenstücken: Der Körper bewegt sich zuerst langsam und gleichförmig und ab dem Knick schneller, wieder gleichförmig. Im Knick ändert sich die Geschwindigkeit sprunghaft.
\end{abcliste}
